\textbf{结论名称：} 球面距离（大圆劣弧长）

\textbf{适用条件：} 半径为 \(R\) 的球面上两点，已知经纬度或球心角 \(\theta\)。

\textbf{变量说明：} 
\(R\)：球半径； 
\(A(\varphi_1,\lambda_1),\; B(\varphi_2,\lambda_2)\)：纬度与经度（弧度），北纬东经为正； 
\(\Delta\lambda = \lambda_2 - \lambda_1\)； 
\(\theta = \angle AOB\) 为球心角。

\textbf{核心公式：}
\[
\boxed{l = R \arccos\!\bigl( \sin\varphi_1\sin\varphi_2 + \cos\varphi_1\cos\varphi_2\cos(\Delta\lambda) \bigr)}
\]
或等价地 \(l = R\theta\)，其中 \(\theta\) 可由球面余弦定理求得。

\textbf{使用场景：} 计算地球表面两点间最短路径（大圆航线），用于航空、航海及卫星覆盖距离估算。

\textbf{简单例子：}
单位球 \(R=1\)，赤道上经度差为 \(90^\circ\) 的两点：
\(A(0,0),\; B(0,\frac{\pi}{2})\)。
代入公式：
\(\sin 0\sin 0 + \cos 0\cos 0\cos\frac{\pi}{2} = 0\)，
\(\arccos 0 = \frac{\pi}{2}\)，
故球面距离 \(l = 1 \times \frac{\pi}{2} = \frac{\pi}{2}\)。